\scanpage{65}
\uline{Supposons} que $\mathbb W^a(\mathbb W(S)[t])$ existe. C'est alors de façon naturelle un préschéma en anneaux sur $S$, que nous noterons $W$, appelé \uline{préschéma d'anneaux générateur du foncteur} $\mathbb W$.

[Dans l'exemple II, c'est le schéma d'anneaux $W$ de départ. Dans le cas de l'exemple I, \deleted{c'est le} si \margin{\unsure{i)|')}} \deleted{schéma} on suppose $S'$ \inserted{fini et} loc. libre sur $S$, c'est le \deleted{\unsure{faisc|faisc.}} schéma en anneaux $V(\check{\underline A})$, où $\underline A$ est le faisceau d'algèbres sur $S$ associé à $S'$].

\margin{\unsure{')|i)}}Soit alors $\unsure{\mathfrak X|\mathcal X}$ ss-schéma fermé de $\mathbb W(S)[t_1,\ldots,t_r]$ défini par \deleted{\unsure{un}} le faisceau d'idéaux dans \margin{\unsure{$\alpha$|a}} $\underline O_{\mathbb W(S)}[t_1,\ldots t_r]$ engendré par des sections $F_i$ $(i=1,\ldots s)$, donc $\mathfrak X$ est le produit fibré, sur $\unsure{\mathfrak Y|\mathcal Y}^s$, de $\mathfrak Y^r$ et $\mathfrak Y^0$ [où $\unsure{\underline{\mathfrak Y}|\mathfrak Y|\mathcal Y}\deleted{^{\unsure{*|1}}}=\mathbb W[S][t]$).
% archive sheet number: 64
